% Section 10.1: 01-charts-and-atlases | Part: 2 | Depends on: part:algebra | Status: done | Original lines: 1805-1849
%!TEX root = ../../../main.tex
\begin{intuition}[{a space that looks flat up close}]
A \emph{manifold} $M$ is a set that, near each point, looks exactly like $\R^m$, even though globally it may be curved, closed, or have holes: a circle looks like a line locally; a sphere or a torus looks like a plane locally; the Earth looks flat to an ant. A manifold is a ``space without measures'': it knows \emph{which} points are near each other and what it means for a function or a curve to be smooth, but it has no notion of distance or angle yet (that is the role of the metric, later). In GR, \textbf{spacetime is a 4-dimensional manifold} whose points are events.

An \textbf{atlas} is a collection of ``maps'' (charts) that together cover $M$, like the pages of an atlas of the world; \textbf{coordinates} are what a chart assigns to each point. Overlapping charts must agree smoothly: the \emph{transition maps} must be $C^\infty$.
\end{intuition}

\begin{formal}[{smooth manifold (via atlases)}]
Let $M$ be a set. An \emph{$m$-dimensional smooth structure} is a collection of pairs $(U_\alpha,\varphi_\alpha)$ (\emph{charts}) with $U_\alpha\subset M$ and $\varphi_\alpha:U_\alpha\to\R^m$ such that
\begin{enumerate}[label=(\roman*),nosep]
\item each $\varphi_\alpha$ is a bijection onto an open set $\varphi_\alpha(U_\alpha)\subset\R^m$;
\item $\varphi_\alpha(U_\alpha\cap U_\beta)$ and $\varphi_\beta(U_\alpha\cap U_\beta)$ are open;
\item (\emph{smooth compatibility}) when $U_\alpha\cap U_\beta\ne\emptyset$ the \emph{transition map}
$\varphi_\beta\circ\varphi_\alpha^{-1}:\varphi_\alpha(U_\alpha\cap U_\beta)\to\varphi_\beta(U_\alpha\cap U_\beta)$ is $C^\infty$;
\item the $U_\alpha$ cover $M$ with countably many sets (``second countable'');
\item any two distinct points are contained in disjoint chart domains, or in a common one (``Hausdorff'').
\end{enumerate}
(iv) and (v) are technical topological conditions. Two atlases whose charts are all compatible define the same smooth structure.

A map $\Phi:M\to N$ between manifolds is \emph{smooth} if its coordinate representation $\psi\circ\Phi\circ\varphi^{-1}$ is $C^\infty$ for all charts. A \emph{diffeomorphism} is a smooth bijection with smooth inverse (then $\dim M=\dim N$). $C^\infty(M)$ is the set of smooth real functions on $M$.
\end{formal}

% REWRITTEN: the angle is written theta (phi is reserved for maps; the charts keep their names varphi_alpha; docs/NOTATION.md)
\begin{example}[{the circle needs two charts}]
$S^1=\{(x,y):x^2+y^2=1\}$. The angle $\theta$ is a good coordinate except at one point, where it jumps. Take
\[
U_1=S^1\setminus\{(-1,0)\},\ \ \varphi_1=\theta\in(-\pi,\pi),\qquad
U_2=S^1\setminus\{(1,0)\},\ \ \varphi_2=\theta\in(0,2\pi).
\]
On the overlap (the circle minus two points) the transition map is $\varphi_2\circ\varphi_1^{-1}(\theta)=\theta$ for $\theta\in(0,\pi)$ and $=\theta+2\pi$ for $\theta\in(-\pi,0)$: a translation, hence smooth. (The notes use the arcs $(-\tfrac34\pi,\tfrac34\pi)$ and $(\tfrac14\pi,\tfrac74\pi)$, same idea.) One chart is impossible: $S^1$ is compact while an open subset of $\R$ is not (a bijective continuous map would have to preserve this).
\end{example}

\begin{example}[{the 2-sphere: stereographic projection}]
Project from the north pole $N=(0,0,1)$ and from the south pole $S=(0,0,-1)$:
\[
\varphi_N(x,y,z)=\Big(\tfrac{x}{1-z},\tfrac{y}{1-z}\Big)=:u,\qquad
\varphi_S(x,y,z)=\Big(\tfrac{x}{1+z},\tfrac{y}{1+z}\Big)=:u'.
\]
Since $x^2+y^2=1-z^2$, one finds $|u|^2=\frac{1+z}{1-z}$ and $u'=u\cdot\frac{1-z}{1+z}=\dfrac{u}{|u|^2}$.
The transition map $u\mapsto u/|u|^2$ is smooth on $\R^2\setminus\{0\}$. So two charts suffice. The notes use two charts of spherical coordinates $(\theta,\varphi)$ with different angular domains to the same effect.
\end{example}

\begin{pitfall}[{non-examples}]
The union of two lines crossing in an ``X'' is \emph{not} a manifold at the crossing (no neighbourhood looks like an interval). The cone $z^2=x^2+y^2$ is not a smooth submanifold at its tip. A manifold must look like $\R^m$ near \emph{every} point.
\end{pitfall}
